\documentclass[11pt,a4paper]{article}
\usepackage[margin=25mm]{geometry}
\usepackage{amsmath,amssymb,amsthm,hyperref}
\setlength{\emergencystretch}{2em}
\newtheorem{lemma}{Lemma}\newtheorem{theorem}[lemma]{Theorem}
\newtheorem{corollary}[lemma]{Corollary}
\newcommand{\gen}[1]{\langle#1\rangle}
\newcommand{\Sym}{\operatorname{Sym}}
\title{Pure cuts and the stable counting correspondence}
\author{Human Compression project}\date{}
\begin{document}\maketitle
All groups in the proof are actual permutation groups. The notation
$\Sym(X)$ means the full symmetric group supported on $X$, fixing its
complement pointwise. A perfect cut of a transitive string C-marking
has two omission orbits $A,B$, with the preceding and following label
intervals supported and transitive on $A,B$, respectively. Its mark is
$(a\ b)\alpha\beta$, where the internal factors are supported on
$A\setminus\{a\}$ and $B\setminus\{b\}$. A pure cut has mark $(a\ b)$.
Internal cuts exclude the two endpoint labels.

We acquire the exact leaf-kernel criterion, the actual-wing bounds and
the complete sufficient-split theorem. The latter gives, when $2r>n+2$,
an internal pure cut or an internal perfect cut with one trivial internal
factor and two odd actual attaching groups. It does not assert that every
usable cut admits subdivision. We shall use pure cuts throughout the
counting argument.

\begin{lemma}[A symmetric factor admitting a pure leaf]\label{backwards}
Let $H=\gen{h_0,\ldots,h_{k-1}}=\Sym(W)$ be a natural symmetric
string C-marking, with $a$ moved by $h_0$ and fixed by the later marks.
If adjoining $(a\ c)$ at a new point gives a C-marking, then $h_0$
is a pure transposition. More generally the same conclusion follows
from the full leaf-kernel tests at $a$.
\end{lemma}
\begin{proof}
Put $P_j=\gen{h_0,\ldots,h_{j-1}}$ and $X_j=a^{P_j}$, with
$P_0=1$. The acquired exact criterion gives
\[
 P_j\cap\Sym(X_{j-1})\le P_{j-1}\qquad(1\le j\le k).
\]
We have $P_k=\Sym(W)\supseteq\Sym(X_k)$. Descending on $j$, if
$P_j$ contains $\Sym(X_j)$, it contains $\Sym(X_{j-1})$, and the
$j$th test puts the latter group in $P_{j-1}$. Hence every $P_j$
contains $\Sym(X_j)$. In particular $\gen{h_0}$ contains the pure
transposition on $X_1=\{a,a^{h_0}\}$. Since it has order two, that
transposition is $h_0$ itself.
\end{proof}

\begin{theorem}[Existence of an internal pure cut]\label{exist}
A fractured natural symmetric string C-marking of degree $n\ge5$ and
rank $r$ with $2r>n+2$ has an internal pure perfect cut. For $n\ge7$
fracture is automatic. In degrees five and six the ordinary simplex
representative gives the same conclusion on marked-group isomorphism
classes.
\end{theorem}
\begin{proof}
Apply the acquired sufficient-split theorem. If its cut is not pure,
reverse the order if necessary and write it as
$q=(a\ b)\beta$, with supported transitive wings $L,R$ on $A,B$.
The attaching groups $L$ and $J=\gen{q,R}$ are both odd. Set
$p=\operatorname{rank}(L)$ and $s=\operatorname{rank}(R)$. Then
\[
                    (2p-|A|)+(2s-|B|)>0.                \tag{1}
\]
If $2s>|B|$, the actual-wing theorem gives $R=\Sym(B)$.
Since $\beta\in R$, it follows that $J=\Sym(\{a\}\cup B)$.
The original IP and the odd external group $L$ give the full
leaf-kernel tests for $J$ at $a$, by the acquired odd-attachment
lemma. Lemma~\ref{backwards} makes its first mark $q$ pure,
contrary to this case.

Thus $2p>|A|$, and $L=\Sym(A)$. Apply the same full-kernel
argument to $L$, oriented from its last mark, with odd external
$J$. Its last mark is pure, say $(d\ a)$. All earlier $L$ marks
fix $a$. Their graph on $A\setminus\{a\}$ is connected: deleting
a pendant vertex from the connected graph of $L$ preserves
connectedness. They are therefore transitive on that set. In the
whole marking the opposite interval is $J$, supported and transitive
on $\{a\}\cup B$. Hence $(d\ a)$ is a pure perfect cut of the
whole marking. It is internal: $2p>|A|$ forces $p\ge2$, since a
single transitive involution has degree two; the original $q$ and
its nonempty suffix follow it.

The acquired initial-fracture theorem supplies fracture for $n\ge7$.
At degrees five and six the strict rank inequality forces the simplex
up to marked-group isomorphism. The outer degree-six image is replaced
by the ordinary action, exactly as in that acquired classification.
\end{proof}

\begin{lemma}[The interval of pure cuts]\label{interval}
In any transitive string C-marking, every label between two pure
perfect cuts is itself a pure perfect cut. Consecutive such cuts
form a simple path, whose internal vertices are fixed by every
mark outside that path.
\end{lemma}
\begin{proof}
Let the cuts have labels $i<j$, with omission supports $A_i,B_i$
and $A_j,B_j$. We have $A_i\subset A_j$ and $B_j\subset B_i$:
the earlier prefix is contained in the later prefix, and the pure
$i$-edge, also in the latter, joins its support to one further point.
This also proves the inclusion when the earlier prefix is trivial
on a singleton. The reversed argument gives the other inclusion.
Put $W=A_j\setminus A_i=B_i\setminus B_j$. Write the two marks
as $u=(a\ b)$ and $v=(c\ d)$, with $a\in A_i$, $b,c\in W$
and $d\in B_j$.

If $j=i+1$, the later prefix is generated by the earlier one and
$u$. Its orbit is exactly $A_i\cup\{b\}$, so $W=\{b\}$ and
$b=c$. Every other mark fixes this point, by the two separate
support conditions.

Suppose $j>i+1$, and let $M$ be the intervening label subgroup.
It is supported on $W$. It is transitive there: the suffix at $i$
is generated by $M,v$ and the suffix at $j$, and its sole edge
from $W$ to the latter's orbit $B_j$ is $cd$. Thus every point
of $W$ lies in the $M$-orbit of $c$. The pure-leaf factorization
now gives the actual supported groups
\[
 \gen{u,M}=\Sym(\{a\}\cup W),\qquad
 \gen{M,v}=\Sym(W\cup\{d\}).
\]
Their original label sets intersect exactly in those of $M$.
The original IP therefore gives $M=\Sym(W)$.

The first mark of $M$ moves $b$, and all its later marks fix $b$
by commutation with $u$. Lemma~\ref{backwards}, applied to the old
C-submarking $(u,M)$, makes that first mark pure. It is a perfect
cut globally: its prefix is the old prefix joined through $u$ to
$b$, and its suffix is the intervening tail on $W\setminus\{b\}$
joined through $v$ to $B_j$. The tail is transitive after deleting
the pendant $b$. Here $b\ne c$, since the distant pure marks
$u,v$ commute. The argument also covers a one-mark intervening
subgroup, whose remaining support is a singleton.

Repeat with this new cut and $j$. This proves the interval assertion.
The adjacent-cut calculation shows that successive middle vertices
are fixed by every other mark. They are distinct: each new prefix
orbit gains exactly its new middle vertex. Thus the cuts are a
simple path with the stated support property.
\end{proof}

\begin{lemma}[The first internal pair]\label{pair}
If $2r>n+2$ and an internal pure perfect cut exists, the least
internal pure index $i$ is followed by another internal pure perfect
cut at $i+1$.
\end{lemma}
\begin{proof}
Let its wing degrees be $a,b$ and ranks $p=i$, $s=r-i-1$.
Then $(2p-a)+(2s-b)>0$. If $2p>a$, the actual-wing theorem makes
the prefix $\Sym(A)$. Its old pure-leaf extension by the $i$-mark
is C. Lemma~\ref{backwards} and the pendant argument in
Theorem~\ref{exist} give a pure perfect cut at $i-1$. It is
internal because $p\ge2$, contradicting the choice of $i$.
Hence $2s>b$. The same argument on the symmetric suffix gives
a pure perfect cut at $i+1$. Since $s\ge2$, this cut is internal.
\end{proof}

\begin{lemma}[Contraction of a pure pair]\label{contract}
Two consecutive pure perfect cuts may be contracted to one pure
perfect cut. The result is a string C-marking of the full symmetric
group on one fewer point, and has one fewer mark.
\end{lemma}
\begin{proof}
The adjacent geometry in Lemma~\ref{interval} gives disjoint sets
$A,\{c\},B$, supported transitive preceding and following groups
$L,T$, and the two pure marks $(a\ c),(c\ b)$. Every other mark
fixes $c$. Replace this pair by $q=(a\ b)$ and delete $c$.
The new mark is distinct from all old ones, which are supported
separately on $A,B$. Every old distant mark fixes the corresponding
endpoint by the old string relations, so the new string relations hold.

We pay IP through actual intersections. For subgroups $P\le\Sym(A)$
and $Q\le\Sym(B)$ put $X=a^P$, $Y=b^Q$ and
$H(P,Q)=\gen{P,Q,(a\ b)}$. Conjugation of $(a\ b)$ by $P,Q$
gives every $(x\ y)$ with $x\in X,y\in Y$. These generate the
symmetric group on $X\cup Y$. Cancelling restrictions on that set
gives the direct-product identity
\[
 H(P,Q)=\Sym(X\cup Y)\times P^{A\setminus X}
                         \times Q^{B\setminus Y}.        \tag{2}
\]
The superscripts denote actual restriction images, supported on
the displayed complements. Consequently, whenever $P_0\le P_1$
and $Q_0\le Q_1$,
\[
 H(P_0,Q_1)\cap H(P_1,Q_0)=H(P_0,Q_0).                  \tag{3}
\]
Indeed an intersection element preserves the intersection
$X_0\cup Y_0$ of the two large symmetric supports. Its complementary
actions belong to the images of $P_0$ and $Q_0$, respectively.
These three sets partition the domain, so (2) proves its membership
in the right-hand side. The reverse inclusion is immediate.

Every interval ending at $q$ is a point renaming of an old interval
ending at $(a\ c)$: rename its point $c$ as $b$. Every interval
beginning at $q$ comes similarly from the old interval beginning at
$(c\ b)$. These and intervals not containing $q$ are C. Induct
on the length of an interval containing $q$ and marks on both sides.
Its end deletions are C by induction. Their groups are the two groups
on the left of (3), and their common middle is the right-hand side.
The acquired standard end criterion proves that interval C, and
finally proves the whole marking C.

For the full transitive $L,T$, formula (2) is $\Sym(A\cup B)$.
Its omission at $q$ is the supported product of those two groups,
so $q$ remains a pure perfect cut. Subdivision of $q$ recovers the
original marking literally after naming the new point $c$.
\end{proof}

\begin{lemma}[Transport of pure cuts]\label{transport}
Subdivision of a specified pure perfect cut at $i$ replaces its
pure-cut index by $i,i+1$. All other pure perfect cuts are retained,
with the later indices shifted by one, and there are no further
pure perfect cuts. Contraction in Lemma~\ref{contract} reverses this
bookkeeping.
\end{lemma}
\begin{proof}
Write the old cut $q=(a\ b)$ and its new marks $u=(a\ c)$,
$v=(c\ b)=uqu$, with $c$ new. For an old index $j\ne i$,
the new maximal omission is literally $\gen{G_j,u}$, because
$q=uvu$. Its orbits are the old orbits, with $c$ adjoined to the
one containing $a$: adjoining a transposition at a new point joins
exactly that point and that one old orbit.

For $j<i$, an old perfect cut has $a$ in its right orbit, since
its suffix contains $q$ and is supported there. The new prefix is
unchanged and its suffix is the old suffix enlarged by $u$. Thus
the cut remains perfect, with the same crossing mark and internal
factors. Conversely a new perfect cut before $i$ has $a,c$ in its
right orbit, since its suffix contains $u$. The orbit identity
shows that deleting $c$ gives precisely the old supported transitive
suffix. The old prefix and crossing mark are unchanged. This proves
the converse as well, including purity. Reverse the argument for
$j>i$. The two inserted cuts are pure perfect cuts directly from
their supported wings. Contraction is the literal inverse
construction of Lemma~\ref{contract}, so has the reverse index rule.
\end{proof}

\begin{theorem}[Stable counting]\label{stable}
Let $\Sigma^\kappa(n)$ be the rank-$n-\kappa$ string C-markings
of $S_n$, modulo marked-group isomorphism and reversal of all labels.
For every $\kappa\ge1$ its cardinality is constant for
$n\ge2\kappa+3$.
\end{theorem}
\begin{proof}
First use ordered point-conjugacy classes and $n\ge7$. Put
$r=n-\kappa$. The domain inequality is $2r\ge n+3$.
Theorem~\ref{exist} supplies an internal pure perfect cut. Define
$E$ by subdividing the least such index. The acquired pure-split
theorem proves the new marking C. It generates $S_{n+1}$ because
it contains the old $S_n$ fixing the new point and a transposition
joining that point to the old domain. Its rank is $r+1$.

Conversely a target of degree $N=n+1$ and rank $R=r+1$ satisfies
$2R\ge N+4$. Theorem~\ref{exist} and Lemma~\ref{pair} supply
two internal pure perfect cuts starting at its least internal pure
index. Define $D$ by their contraction. Lemma~\ref{contract}
proves actual IP and full symmetric generation. The contracted
parameters satisfy $2(R-1)\ge(N-1)+3$, the exact domain condition.
Its retained cut stays internal because both original cuts were
internal. Lemma~\ref{transport} preserves the least internal index
under either operation. Consequently, as ordered markings up to
renaming the inserted point,
\[
                         D(E(X))=X,\qquad E(D(Y))=Y.     \tag{4}
\]

We must also pay reversal. By Lemma~\ref{interval}, all pure perfect
cuts belong to one consecutive path. Subdividing any internal edge
of this path gives the same longer path with the same endpoints,
in its final ordered labels. All marks outside the interval have
the same positions and actions. The explicit conjugacy fixes every
point outside the internal path vertices, and matches the vertices
in order from the same first endpoint. It intertwines each marked
permutation. Hence extension at any two internal pure cuts gives
the same ordered isomorphism class. Reversal reflects the interval
and sends subdivision at $i$ to subdivision at $r-1-i$ in the
reversed old marking. The independence of the pure choice therefore
makes $E$ commute with reversal on ordered classes. Its inverse
$D$ does too. Equations (4) descend to the quotient by reversal,
including self-dual classes.

For completeness, ordered marked-group isomorphism is point conjugacy
for $S_n$ when $n=5$ or $n\ge7$. Among its nonidentity involution
classes the transpositions have uniquely smallest size. For a product
of $k\ge2$ disjoint transpositions, the ratio of its class size to
the transposition class size is
\[
                  \frac{(n-2)!}{2^{k-1}k!(n-2k)!}>1.
\]
For $k=2$ this follows immediately for $n\ge5$; for $k=3$ use
$n\ge7$; for $k\ge4$ the minimum occurs at $n=2k$, starts at
$6!/(8\cdot4!)>1$, and increases with $k$. Thus every automorphism
preserves transpositions. The maximal pairwise noncommuting families
of at least four transpositions are exactly the stars at one point:
a pairwise intersecting family of edges is a star or a triangle.
Automorphisms therefore permute the point stars, and their pairwise
intersections identify each transposition. Since these generate $S_n$,
the automorphism is conjugation by that point permutation.

The only transitions involving degree five or six in the stated range
have $\kappa=1$. Their acquired complete extremal classification is
the simplex up to marked-group isomorphism. Choose its ordinary action;
the outer degree-six image is the same marked-group class. Subdivision
and contraction give the successive simplex classes, so (4) and reversal
hold there as well. This pays all isomorphism conventions.

We have proved $|\Sigma^\kappa(n)|=|\Sigma^\kappa(n+1)|$ at every
$n\ge2\kappa+3$. Iteration gives the assertion.
\end{proof}

\paragraph{Attribution and scope.}
The counting theorem and the rank and degree operation are due to
Cameron, Fernandes and Leemans, \emph{The number of string C-groups of
high rank}, arXiv:2212.12723v2. The acquired classifications and finite
boundaries retain their original attributions. The present argument
uses the repaired leaf criterion, the completed sufficient-split proof,
backward propagation of its full kernels, and the interval of pure cuts.
It pays both inverse equations and the quotient by reversal without
asserting that the printed unrestricted extension criterion is valid.
The standard end criterion and supported-transposition factorization
are acquired with their earlier proofs. No numerical stable constants
are computed here; their exact base classifications remain separate.
This manuscript is a proof component pending distinct review, with no
whole-original, new-conjecture, Lean or compression credit assigned.
\end{document}
